\documentclass[preprint,10pt]{elsarticle}
% \documentclass[3p,twocolumn]{elsarticle}
\usepackage[bookmarks=true,colorlinks=true,linkcolor=blue]{hyperref}

\usepackage{amsmath,amssymb,mathabx}
\usepackage{graphicx,esint,ulem}

\biboptions{sort&compress}% compress bibtex entries
\usepackage[usenames,dvipsnames,svgnames,table]{xcolor}

\usepackage{color}
\definecolor{jwbGreen}{rgb}{0, .6, 0}
\definecolor{jbaPurple}{HTML}{6600FF}
\definecolor{purple}{rgb}{.7, 0., .8}
\newcommand{\red}{\color{red}}
\newcommand{\green}{\color{green}}
\newcommand{\blue}{\color{blue}}
\newcommand{\purple}{\color{purple}}
\newcommand{\orange}{\color{orange}}

\usepackage{calc}

% for fancy coloured boxes: 
\usepackage{empheq}
\usepackage[most]{tcolorbox}

\newtcbox{\mymath}[1][]{%
    nobeforeafter, math upper, tcbox raise base,
    enhanced, colframe=blue!60!black,
    colback=blue!20, boxrule=1pt,
    #1}


\usepackage{algorithm,algpseudocode}
\usepackage{algorithmicx}
\algrenewcommand\alglinenumber[1]{\footnotesize #1:} % Algorithm line number font size
\newcommand{\algFontSize}{\footnotesize}

\usepackage{tcolorbox}
\usepackage[english]{babel}
\usepackage{amsmath}
\usepackage{amsfonts}
\usepackage{fullpage}
\usepackage{graphicx}
\usepackage{setspace}
\usepackage{float}

\usepackage{tabu}% to increase table row separation

%\usepackage{showkeys}
\usepackage{multirow}

\usepackage{fancyvrb}% fancy verbatim


% \usepackage{mdframed}% for shaded framed boxes


% ************ ALLOW displayed equations to cross pages *********************
\allowdisplaybreaks

% ---------------------------- DATE HEADER ------------------------------------
% \usepackage{fancyhdr}
% \usepackage{time}
% take the % awaNOy on next line to produce the final camera-ready version
% Be sure to remove \thispagestyle{fancy} as well after the \make.
% \pagestyle{empty}
%
%\pagestyle{fancy}
% \newcommand\myTime{\now}
%\fancyhead{}
%\fancyhead[CO, CE]{\texttt{-Draft-}}
%\fancyhead[RO, RE]{\texttt{\today, \myTime}}
%\setlength{\headheight}{2\baselineskip}
%\renewcommand{\headrulewidth}{0pt}
% ---------------------------- DATE HEADER ------------------------------------



\usepackage{tikz}

\newtheorem{theorem}{Theorem}
\newtheorem{assumption}{Assumption}
\newtheorem{definition}{Definition}
\newenvironment{proof}[1][Proof]{\begin{trivlist}
\item[\hskip \labelsep {\bfseries #1.}]}{\end{trivlist}}

\newtheorem{thm}{Theorem}[section]
\newtheorem{prop}{Proposition}[section]
\newtheorem{lem}{Lemma}[section]
%\newtheorem{cor}{Corollary}[section]
\newtheorem{example}{Example}[section]
\newtheorem{conj}{Conjecture}[section]
\newtheorem{remark}{Remark}[section]
% some definitions of bold math italics to make typing easier.
% They are used in the corollary.

 
\input tex/defs.tex
\input tex/trimFig.tex
\usepackage{xargs}% for optional args to \newcommandx
\input tex/plotFigureMacros.tex

% \newcommand{\new}[1]{{\color{olive}#1}}
\newcommand{\new}[1]{{\color{red}#1}}
\newcommand{\wdh}[1]{{\color{DarkBlue}wdh: #1}}
\newcommand{\old}[1]{{\color{red}old: #1}}
\newcommand{\newer}[1]{{\color{orange} #1}}
\newcommand{\avk}[1]{{\color{orange}AVK: #1}}

\usepackage[margin=.75in]{geometry}

\renewcommand\floatpagefraction{.99}
\renewcommand\topfraction{.99}
\renewcommand\bottomfraction{.99}
\renewcommand\textfraction{.01}   
\setcounter{totalnumber}{50}
\setcounter{topnumber}{50}
\setcounter{bottomnumber}{50}
% ============================================================================================
\begin{document}



\begin{frontmatter}
 \title{Documentation for the incompressible Navier-Stokes Matlab codes}


% \author[rpi]{Cassandra~Carrick\fnref{NSFgrants}}
% \ead{carric2@rpi.edu}


% \author[rpi]{Jeffrey W.~Banks\fnref{DOE}}
% \ead{banksj3@rpi.edu}

\author[rpi]{William D.~Henshaw\corref{cor}\fnref{NSFgrants}}
\ead{henshw@rpi.edu}

% \author[rpi]{Donald~W.~Schwendeman\fnref{NSFgrants}}
% \ead{schwed@rpi.edu}


\address[rpi]{Department of Mathematical Sciences, Rensselaer Polytechnic Institute, Troy, NY 12180, USA}

% \cortext[cor]{Corresponding author}

% % \fntext[DARPA]{This work was partially funded by the DARPA Defense Sciences Office, Award HR00111720032.}
% \fntext[NSFgrants]{Research supported by the National Science Foundation under grants DMS-1519934 and DMS-1818926.}


% \fntext[RTG]{This work was partially funded by the NSF Research Training Group Grant DMS-1344962.}
% \fntext[DOE]{This work was partially performed under DOE contracts from the ASCR Applied Math Program.}
% %\fntext[DOD]{This work was partially performed under DOD contract W911NF-14-C-0161.}

% % \fntext[NSFgrantNew]{Research supported by the National Science Foundation under grant DMS-1519934.}

% \fntext[PECASEThanks]{Research supported by a U.S. Presidential Early Career Award for Scientists and Engineers.}

\begin{abstract}

  Here is documentation and notes for the incompressible Navier-Stokes Matlab codes.


\end{abstract}

\end{frontmatter}


\tableofcontents

% \clearpage
% \input tex/damping 


\newcommand{\Vdot}{\dot{V}}
\newcommand{\lambdaHat}{\hat{\lambda}}


% -----------------------------------------------------------------
\clearpage
\section{Matlab files}

\begin{table}[H]\small % you should set \tableFont to \footnotesize or other size
% \newcommand{\num}[2]{#1e{#2}} % use this command to set the format of numbers in the table.
\begin{center}
\begin{tabular}{ll} % \hline 
  ins.m                              &  Main routine . \\
  runChecks.m                        &  Run regression tests. \\
  runGridConvergence.m               &  Run a grid convergence test. \\
  advanceAdams.m                    &  Adams time-stepping. \\
  advanceIM.m                       &  IMEX time-stepping. \\
  advancePC.m                       &  Predictor-Corrector time-stepping. \\
  applyBoundaryConditions.m         & apply (explicit) boundary conditions  \\
  pressureEquation.m                & form pressure equation and solve \\
  evalMap.m                         & eval the grid mapping and derivatives \\
  formImplicitTimeSteppingMatrix.m  & form systems for implicit time-stepping \\
  solveImplicitTimeStep.m           & solve the implicit time-step equations \\
  getInitialConditions.m            & assign initial conditions \\
  getUt.m                           & evaluate momentum equations \\
  getTimeStep.m                     & determine the time-step $\dt$ \\
  plotSolution.m                    & plot solution contours and streamlines \\
  plotStreamLines.m                 & stream line plotter 
\end{tabular}
  \caption{Selected ins Matlab files.}
\end{center}
\end{table}


\input tex/intro

\input tex/discreteApproximations

% % ------------------------------------------------------------------------------
% \section{Grid convergence tests}

% The Matlab code \texttt{runGridConvergence.m} can be used to run grid convergence tests.
% Here is an example
% \begin{Verbatim}[fontsize=\small]
% >> runGridConvergence -ts=pc2 -tf=.25 -ms=trig -numResolutions=4
%   N     p-err  ratio   u-err  ratio   v-err ratio    div  ratio    cpu(s)  ratio
%  --------------------------------------------------------------------------------- 
%   10  4.46e-02       6.31e-03       4.17e-03       4.00e-02        1.00e-02
%   20  1.16e-02  3.8  4.71e-04 13.4  3.54e-04 11.8  3.28e-03 12.2   8.00e-02  8.0
%   40  2.92e-03  4.0  1.17e-04  4.0  6.59e-05  5.4  2.31e-04 14.2   1.35e+00 16.9
%   80  7.41e-04  3.9  2.92e-05  4.0  1.57e-05  4.2  1.79e-05 12.9   1.01e+01  7.5
% \end{Verbatim}



\input tex/numericalResults




\clearpage
\bibliographystyle{elsart-num}
\bibliography{journal-ISI,ins}
 
\end{document}


% ***************************************************************************************
% ***************************************************************************************
% ***************************************************************************************
% ***************************************************************************************
% ***************************************************************************************
% ***************************************************************************************
% ***************************************************************************************









\begin{table}[H]\footnotesize % you should set \tableFont to \footnotesize or other size
% \newcommand{\num}[2]{#1e{#2}} % use this command to set the format of numbers in the table.
\begin{flushleft}
\begin{tabular}{ll} % \hline 
   % % ************ START OPTIONS FOR COMPUTE EIGS **********
   %  par.evFileName              = getString( line,'-evFileName',par.evFileName );
   %  par.plotEvects              = getInt( line,'-plotEvects',par.plotEvects );
   %  par.upwindScaleFactor       = getReal( line,'-upwindScaleFactor',par.upwindScaleFactor ); 
   %  % ************ END OPTIONS FOR COMPUTE EIGS **********    
     % for( ipc=1:pcMaxLevels ) 
    %   par.pc{ipc}.its         = getInt( line,sprintf('-pc%dIts',ipc),par.pc{ipc}.its );
    %   par.pc{ipc}.damp        = getReal( line,sprintf('-pc%dDamp',ipc),par.pc{ipc}.damp );  
    % end  

    ms                      & [poly$\vert$pulse$\vert$gaussian$\vert$eig] : manufactured solution or solution type\\
    useWaveHoltz            & [0$\vert$1] :  0=solve wave equation, 1=solve waveHoltz or EigenWave\\
    checkFileName           & name of the check file for regression tests.\\
    help                    & print help information.\\
    debug                   & debug bit flag.\\
    nd                      & number of dimenions (1 or 2)\\
    order                   & order of accuracy (2 or 4)\\
    numGhost                & \\
    plotOption              & \\
    pauseOption             & \\
    numResolutions              & \\
    N0                          & \\
    implicit                & \\
    implicitFirstStep       & \\
    filterTimeDerivative    & \\
    adjustOmega             & \\
    useDiscreteBetaFunction & \\
    useFilterWeights        & \\
    usePeriodicFirstStep    & \\
    checkForcingOnBoundary  & \\
    formWaveHoltzMatrix     & \\
    readWaveHoltzMatrix     & \\
    coarseGridPreconditioner& \\
    numPC                   & \\
    numFixPointIterations       &  maximum number of fixed point iterations \\
    par.maxit                   &  maximum number of Krylov iterations \\
    par.restart                 & \\
    numFreq                     & \\
    par.minStepsPerPeriod       & minimum number of WaveHoltz steps per period\\
    par.maxRitz                 & \\
    par.maxEigVectsToPlot       & \\
    numPer1                     & \\
    numPer1                     & \\
    par.useFixPoint             & solve with fixed point iterations (or power method for EigenWave). \\
    par.useGMRES                & solve WaveHoltz with a Krylov method  (or using Matlab eigs for EigenWave).\\
    par.krylovType              & \\
    par.implicitUpwind          & \\
    par.upwind                  & \\
    par.upwindh                 & \\
    par.useUpwindCoeff          & \\
    par.upwindOperator          & \\
    omega1                      & \\
    omega1                      & \\
    omega2                      & \\
    omega3                      & \\
    par.yMin                    & \\
    par.yMax                    & \\
    par.uMin                    & \\
    par.uMax                    & \\
    kx1                         & \\
    kx2                         & \\
    kx3                         & \\
    ax                          & \\
    bx                          & \\
    ay                          & \\
    by                          & \\
    c                           & \\
    cfl                         & \\
    par.tFinal                  & \\
    par.tPlot                   & \\
    alphaImp                    & \\
    betaBar                     & \\
    tol                         & \\
    lambdaTol                   & \\
    par.lambdaMax               & \\
    par.savePlots               & \\
    par.plotName                & \\
    par.figDir                  & \\
\end{tabular}
  \caption{Selected WaveHoltz options.}
\end{flushleft}
\end{table}

\begin{table}[H]\footnotesize % you should set \tableFont to \footnotesize or other size
% \newcommand{\num}[2]{#1e{#2}} % use this command to set the format of numbers in the table.
\begin{flushleft}
\begin{tabular}{ll} % \hline    
    par.deflateWaveHoltz        & \\
    par.deflateWaveHoltz        & \\
    par.deflateForcing          & \\
    par.degreex                 & \\
    par.degreet                 & \\ 
    par.numToDeflate            & \\
    par.deflateToRate           & \\
    par.acg                     & use augmented CG \\
    par.agmres                  & use augmented GMRES \\
    par.abicgstab               & use augmented Bi-CG-STAB \\
    par.useApproxEigs           & \\
    par.forcingScaleFactor      & \\
    par.phase                   & \\
    par.damp                    & \\
    par.isign                   & \\
    par.computeRitz             & \\
    par.deflateRitz             & \\
    par.ritzConvergenceTol      & \\
    par.numEigsToCompute        & \\
    par.removeConstraints       &  remove constraint equations (e.g. boundary conditions) from WaveHoltz or EigenWave iterations\\
    par.beta(1)                 & \\
    par.beta(1)                 & \\
    par.beta(2)                 & \\
    par.beta(3)                 & \\
    par.xg(1)                   & \\
    par.xg(1)                   & \\
    par.xg(2)                   & \\
    par.xg(3)                   & \\
    par.yg(1)                   & \\
    par.yg(1)                   & \\
    par.yg(2)                   & \\
    par.yg(3)                   & \\
    par.amp(1)                  & \\
    par.amp(2)                  & \\
    par.amp(3)                  & \\
    par.legendLoc               & \\
    bc                          & \\
    par.rbc                     & \\
    par.useTwoLevelRBC          & \\
    par.rbcImplicitOption       & \\
    par.superGrid               & \\
    par.superGridh              & \\
    par.superGridWidth          & \\
\end{tabular}
  \caption{Selected WaveHoltz options continued.}
\end{flushleft}
\end{table}


% \begin{description}
%   \item[waveHoltz.m] Main routine for WaveHoltz and EigenWave.
%   \item[computeEigs.m] Routine to compute eigenpairs directly.
%     - advance.m                  : solve the wave equation
%     - applyBoundaryConditions.m  : apply BC's 
%     - deflationSetup.m           : choose eigenvectors to deflate
%     - getEigenvalues.m           : form matrix and compute discrete eigenpairs
%     - getTrueEigenValues.m       : compute true (discrete) eigenvalues
%     - takeTimeStep.m             : 
%     - formImplicit.m             : form the implicit time-stepping matrix
%     - setupGrid.m                : 
%     - getHelmholtzSolution.m     : 
%     - gmresFunction.m            : matrix-vector multiply routine to call advance for WaveHoltz or EigenWave
    
%     - sigma = getIntegrationWeights( Nt,numFreq,dt,Tv )

%     - solveHelmholtz.m           : main routine to solve the Helmholtz equation using WaveHoltz  
% \end{description}




% -----------------------------------------------------------------

\section{Discrete approximation}

Here is the second-order accurate in time scheme 
\bse
\label{eq:MEscheme}
\ba
   \Dpt\Dmt W_\jv^n = L_h\Big[ \alpha_2 W_\jv^{n+1} + \beta_2 W_\jv^{n} +\alpha_2 W_\jv^{n-1} \Big] \qquad \jv \in\Omega_{h,a},  
\ea
($\beta_2=1-2\alpha_2$, and $\alpha_2=0$ for the explicit scheme)
with initial conditions
\bat
   & W_\jv^0 = V_\jv,               \quad&& \jv \in\Omega_{h,a}   \label{eq:MEschemeBC1} \\
   & \Dzt W_\jv^0 = \Vdot_\jv,      \quad&& \jv \in\Omega_{h,a}   \label{eq:MEschemeBC2}
\eat
\ese
where $\Omega_{h,a}$ denotes the active points. 

\noindent
\textbf{Note: In general (for WaveHoltz and EigenWave) we assume that $V_\jv$, and $\Vdot_\jv$ are only known at active points (not, for example, at ghost points)}.

We also need boundary conditions which can be of Dirichlet, Neumann, or mixed type.
Radiation (impedance) boundary conditions will also be used.

% -----------------------------------------------------------------
\input tex/radiationBoundaryConditions

% -----------------------------------------------------------------
\clearpage
\section{First time step} \label{sec:firstStep}

\paragraph{Explicit time-stepping}
Solving for $W_\jv^{-1}$ from~\cref{eq:MEschemeBC2} gives
\ba
   W_\jv^{-1} = W_\jv^1 - 2\dt \Vdot_\jv
\ea
and substituting this into the interior scheme leads to the formula
for the first time-step,
\ba
   W_\jv^{1} = V_\jv + \dt \Vdot_\jv + \f{\dt^2}{2} ( L_h V_\jv  + F_\jv^0 )  .
   \label{eq:firstStepExplciit}
\ea
This is just the same as a Taylor series in time approximation.


\paragraph{Implicit time-stepping}
The implicit interior scheme is, 
\ba
   W_\jv^{n+1} -2 W_\jv^{n} + W_\jv^{n-1}= \dt^2 L_h\Big[ \alpha_2 W_\jv^{n+1} + \beta_2 V_\jv^n +\alpha_2 W_\jv^{n-1} \Big]  + \dt^2 F_\jv^n 
\ea
or
\ba
  [ I -\alpha_2 \dt^2 L_h] W_\jv^{n+1} = 2\Big[ I + \half\beta \dt^2 L_h\Big] W_\jv^n
                                          -\Big[ I - \alpha_2 \dt^2 L_h\Big] W_\jv^{n-1}  + \dt^2 F_\jv^n 
\ea
Substituting $W_\jv^{-1} = W_\jv^1 - 2\dt \Vdot_\jv$
leads to (for $n=0$)
\ba
  [ I -\alpha_2 \dt^2 L_h] W_\jv^{n+1} = \Big[ I + \half\beta \dt^2 L_h\Big] W_\jv^n
                                       + \Big[ I - \alpha_2 \dt^2 L_h\Big] \Vdot_\jv  + \half\dt^2 F_\jv^n 
\ea
Note that the matrix on the LHS is the same as for a general step.

% With Dirichlet, Neumann, or mixed boundary conditions, the boundary (for Dirichlet) or ghost point values for $V_\jv$ can be filled in and then~\cref{eq:firstStep} can be used for
% all active points.

% With Radiation BC's it may be more involved to assign the ghost point values. This is discussed further in Section~\cref{sec:firstStepWithRaditionBCs}.

% -----------------------------------------------------------------
\clearpage
\section{First time step with radiation boundary conditions} \label{sec:firstStepWithRaditionBCs}

Here we describe how to perform the first step with radiation BCs. The primary issue is how to fill in the ghost values at $t=0$.

The first time-step is given by~\cref{eq:firstStep}.  To apply this first-step on the boundary point $j$ where an RBC is applied we may need 
to perform some manipulations.

\mni
\textbf{Case 1.}
The Sommerfeld condition SRBC-O2-E 
\ba
   \Dzt W_\jv^n - \sigma c \Dzx W_\jv^n = 0,   \qquad \jv\in\p\Omega_{h,r}
\ea
can be combined with the second IC~\cref{eq:MEschemeBC2} to give the mixed type BC:
\ba
   \Vdot_\jv- \sigma c \Dzx V_\jv = 0,   \qquad \jv\in\p\Omega_{h,r}
\ea
which can be used to fill in the ghost values for $V_\jv$ and we are good to go.


\mni
\textbf{Case 2.}
The SRBC-O2-IL2 RBC is 
\ba
   & \Dpt W_\jv^n - \sigma c \Dzx (\half W_\jv^{n} + \half W_\jv^{n+1} )= 0,  \qquad \jv\in\p\Omega_{h,r}
\ea
which for $n=0$ is
\ba
   W_\jv^1 = V_\jv + (2\dt) \sigma c \Dzx (\half V_\jv + \half W_\jv^{1} )  \qquad \jv\in\p\Omega_{h,r}
\ea
We also have for explicit time-stepping 
\ba
   W_\jv^{1} = V_\jv + \dt \Vdot_\jv + \f{\dt^2}{2} L_h\Big[  V_\jv  \Big]   \qquad \jv\in\p\Omega_{h,r} 
\ea

\mni
{\blue If we knew $V_{-1}= W_{-1}^0$ then we could advance the solution. It does not seem possible to determine $V_{-1}$ 
with the given conditions. Maybe the ghost point at $j=-1$ needs to be an active point for this RBC? 
This seems consistent with the discussion of eigenvalue problems in Section~\cref{sec:notesOnEigenWaveForCoplexProblems}.
}

% \bigskip
% which can also be written as (replace $n$ with $n-1$)
% \ba
%    & \Dmt W_\jv^n - \sigma c \Dzx (\half W_\jv^{n} + \half W_\jv^{n-1} )= 0,  \qquad \jv\in\p\Omega_{h,r}
% \ea
% Averaging these last two equations gives
% \ba
%    & \Dzt W_\jv^n - \sigma c \Dzx ( \f{1}{4} W_\jv^{n-1} + \half W_\jv^{n} + \f{1}{4} W_\jv^{n+1} )= 0,  \qquad \jv\in\p\Omega_{h,r}
% \ea
% Using~\cref{eq:MEschemeBC2} then gives
% \ba
%    & \Vdot_\jv - \sigma c \Dzx ( \f{1}{4} W_\jv^{-1} + \half W_\jv^{0} + \f{1}{4} W_\jv^{1} )= 0,  \qquad \jv\in\p\Omega_{h,r}, 
% \ea


% We also know
% \ba
%    \Dzt W_\jv^0 = \Vdot_\jv  \qquad \jv\in\p\Omega_{h,r}
% \ea
% as well as
% \ba
%    W_\jv^{\pm 1} = V_\jv + (\pm\dt) \Vdot_\jv + \f{\dt^2}{2} L_h\Big[ \alpha_2 W_\jv^{n+1} + \beta_2 V_\jv +\alpha_2 W_\jv^{n-1} \Big] 
% \ea


% -----------------------------------------------------------------
\clearpage
\section{Notes on EigenWave for complex valued problems} \label{sec:notesOnEigenWaveForCoplexProblems}

\textbf{Continuous eigenvalue problem.}
Eigenwave solves the unforced wave equation,
\bse
\label{eq:waveEquationIBVPforEigenWave}
\bat
 & \p_t^2 w = L w ,           \quad && \xv\in\Omega, ~t>0,  \\
 & w(\xv,0) = v(\xv),           \quad && \xv\in\Omega \\
 & \p_t w(\xv,0) = \vDot(\xv),           \quad && \xv\in\Omega \\
 & \p_t w - c \p_n w = 0,   \quad && \xv\in\p\Omega_r
\eat
\ese
To obtain the correspoding eigenvalue problem substitute
\ba  
     w = e^{\mu t} \phi(\xv) =  e^{i \sigma \lambda t} \phi(\xv)
\ea
(where $\sigma=\pm$ depending on the sign convention for time-periodic solutions),
 to give the quadratic eigenvalue problem
\bse
\label{eq:contintuousEigenvalueProblem}
\ba  
 &  \mu^2\phi = L \phi ,           \quad && \xv\in\Omega \\
 &  \mu \phi - c \p_n \phi = 0,   \quad && \xv\in\p\Omega_r
\ea
or
\ba  
 & -\lambda^2 \phi = L \phi ,           \quad && \xv\in\Omega \\
 & i \sigma \lambda\phi - c \p_n \phi = 0,   \quad && \xv\in\p\Omega_r
\ea
\ese
\textbf{Note:} In waveHoltz.m $\sigma={\rm par.isign}$, with $\sigma=-1$ by default to match the $e^{-i\omega t}$ convention in electromagnetics.

\bigskip
\noindent
\textbf{Eigenfunction expansions.} When $w(\xv,t)$, $v(\xv)$ and $\vDot(\xv)$ are all real valued,
eigenfunction expansions for~\cref{eq:waveEquationIBVPforEigenWave} take the form
\ba
   w(\xv,t) = \sum_{m} c_m e^{i\sigma \lambda_m t} \phi_m(\xv) + \cBar_m e^{-i\sigma \lambdaBar_m t} \phiBar_m(\xv)
\ea
Recall that if $(\lambda_m,\phi_m)$ is an eigenpair then so is $(-\lambdaBar_m,\phiBar_m)$.
Thus we make use of the double the number of eigenfunctions from the quadratic eigenvalue problem.
The real and imaginary parts of the complex coefficient $c_m = a_m + i b_m$ are determined from two the 
initial conditions:
\bse
\ba
  v(\xv)  & = \sum_{m} c_m  \phi_m(\xv) + \cBar_m  \phiBar_m(\xv) 
      =  \sum_{m} a_m  (\phi_m(\xv) + \phiBar_m(\xv) ) + i b_m (\phi_m(\xv) - \phiBar_m(\xv) )  , \\
  \vDot(\xv) & = \sum_{m} c_m i\sigma \lambda_m \phi_m(\xv) - \cBar_m i\sigma \lambdaBar_m \phiBar_m(\xv)
\ea
\ese

\bigskip
\bigskip
\noindent
\textbf{EigenWave solution.} When using the EigenWave FPI (i.e. power iteration) the solution $w(\xv,t)$ 
will generally converge to a single term in the eigenfunction expansion. This takes the form (after scaling by $c_m$)
\ba
   w(\xv,t) &= e^{i\sigma \lambda_m t} \phi_m(\xv) + \f{\cBar_m}{c_m} e^{-i\sigma \lambdaBar_m t} \phiBar_m(\xv), \\
    &=                      \big(\cos(\lambda_m t) \phi_m + \cos(\lambdaBar t)\phiBar_m\big)
       + i \f{\cBar_m}{c_m} \big(\sin(\lambda_m t) \phi_m - \sin(\lambdaBar t)\phiBar_m\big)  ,
\ea
This implies $v^k$ and $\vDot^k$ converge to
\ba
    & v_k = w(\xv,0) = \big( \phi_m +\phiBar_m\big) ,\\
    & \vDot^k = \p_t w(\xv,0) = i \f{\cBar_m}{c_m} \big(\lambda_m \phi_m - \lambdaBar\phiBar_m\big)  
\ea
\textbf{Question:} What determines the phase variable $\cBar_m/c_m$ ? 

\mni
The radiation BC is $\Bc w = w_t - c \p_n w_x = 0 $
for $\xv\in\p\Omega_r$. Note that 
\ba
 \Bc w =   i\sigma \lambda_m e^{i\sigma \lambda_m t} \phi_m(\xv) -i\sigma \lambdaBar_m \f{\cBar_m}{c_m} e^{-i\sigma \lambdaBar_m t} \phiBar_m(\xv)
   -c \Big[ e^{i\sigma \lambda_m t} \p_n\phi_m(\xv) + \f{\cBar_m}{c_m} e^{-i\sigma \lambdaBar_m t} \p_n\phiBar_m(\xv) \Big] 
\ea
Using $c\p_n\phi_m = i\sigma\lambda_m \phi_m$ for $\xv\in\p\Omega_r$ leads to 
\ba
  \Bc w &= i\sigma \lambda_m e^{i\sigma \lambda_m t} \phi_m(\xv) -i\sigma \lambdaBar_m \f{\cBar_m}{c_m} e^{-i\sigma \lambdaBar_m t} \phiBar_m(\xv)
   - \Big[ e^{i\sigma \lambda_m t} i\sigma\lambda_m \phi_m(\xv) + \f{\cBar_m}{c_m} e^{-i\sigma \lambdaBar_m t} i\sigma\lambda_m \phiBar_m(\xv) \Big]  , \\
      & = 0 
\ea
Thus $\Bc w=0$ for $\xv\in\p\Omega_r$ and the radiation BC is satisfied. 




\bigskip
\bigskip
\noindent
\textbf{Discrete eigenvalue problem.}
The discrete eigenvalue problem can be found by substituting
\ba
    W_\jv^n = A^n \Phi_\jv = \left(e^{\mu \dt}\right)^n \Phi_\jv = \left(e^{i \sigma \lambda \dt}\right)^n \Phi_\jv, 
\ea
into the discrete scheme~\cref{eq:MEscheme} to give 
\ba
   \f{ A - 2 + A^{-1} }{\dt^2}\Phi_\jv = L_h( \alpha_2 A + \beta_2 + \alpha_2 A^{-1} ) \Phi_\jv , \label{eq:discreteEigenvalue}
\ea
with corresponding boundary conditions. \Cref{eq:discreteEigenvalue} can be written
\bse
 \label{eq:discreteEigenvalueII}
\ba
   L_h \Phi_\jv &= \f{1}{\dt^2} \f{A - 2 + A^{-1}}{\alpha_2 A + \beta_2 + \alpha_2 A^{-1}} \Phi_\jv
                = -\lambdaTilde^2 \Phi_\jv
\ea
\ese
where we identify $\lambdaTilde$ as the eigenvalue,
\ba
     \lambdaTilde^2 \eqdef  - \f{1}{\dt^2} \f{A - 2 + A^{-1}}{\alpha_2 A + \beta_2 + \alpha_2 A^{-1}} .
\ea
Dirichlet, Neumann or mixed boundary conditions lead to constraint equations,
\ba
   \Bc_h( \Phi_\jv)=0,
\ea
that can be used to eliminate the boundary and or ghost points from~\cref{eq:discreteEigenvalueII}.

\mni
Radiation boundary conditions are a bit more complicated.
Consider the Sommerfeld RBC
\ba
   \Dzt W_\jv^n - \sigma c \Dzx( W_\jv^{n+1} + W_\jv^{n-1} )/2 , \qquad \jv\in\p\Omega_{h,r},
\ea
which transforms to 
\ba
  &  \f{A-A^{-1}}{2\dt} \Phi_\jv - \sigma c \f{A+A^{-1}}{2} \Dzx \Phi_\jv, \qquad \jv\in\p\Omega_{h,r},
\ea
or
\ba
   & \f{1}{\dt} \f{A-A^{-1}}{A+A^{-1}} \Phi_\jv - \sigma c  \Dzx \Phi_\jv, \qquad \jv\in\p\Omega_{h,r}, \\
   & i\sigma \lambdaHat \Phi_\jv - \sigma c  \Dzx \Phi_\jv, \qquad \jv\in\p\Omega_{h,r},
\ea
where
\ba
   i\sigma \lambdaHat = \f{1}{\dt} \f{A-A^{-1}}{A+A^{-1}}
\ea
The problem is that the square of $i\sigma \lambdaHat$ is not equal to $-\lambdaTilde^2$ for this RBC.
Thus this RBC will correspond to a different eigenvalue problem than the discrete version of~\cref{eq:contintuousEigenvalueProblem}.

To form the discrete \textbf{linear} eigenvalue problem set
\ba
  \Psi_\jv = A \Phi_\jv
\ea
and then $A$ times~\cref{eq:discreteEigenvalue} becomes
\ba
   & \f{ A^2 - 2 A + I }{\dt^2}\Phi_\jv = L_h( \alpha_2 A^2 + \beta_2 A + \alpha_2  ) \Phi_\jv ,  \\
   &  \Big( A \Psi_\jv - 2 \Psi_\jv + \Phi_\jv\Big)  = \dt^2 L_h\Big( \alpha_2 A\Psi_\jv + \beta_2 \Psi_\jv + \alpha_2  \Phi_\jv \Big), 
\ea
which implies the generalized eigenvalue problem for the vector $[\Phi_\jv,\Psi_\jv]$ (*check me*)
\bse
\label{eq:discreteEigenvalueGeneralized}
\ba
     & \Psi_\jv = A \Phi_\jv, \\
     & (2I+\beta_2\dt^2 L_h) \Psi_\jv + (-I +\dt^2\alpha_2 L_h) \Phi_\jv  = A \Big( I - \alpha_2\dt^2 L_h \Big) \Psi_\jv
\ea
\ese
Multiplying $A$ times the RBC condition gives
\bse
\label{eq:sommerfeldEigenvalueV1}
\ba
  &  (A^2-I) \Phi_\jv - \sigma c\dt (A^2+I) \Dzx \Phi_\jv, \qquad \jv\in\p\Omega_{h,r}, \\
  &  A \Psi_\jv - \Phi_\jv - \dt \sigma c \Big( A \Dzx \Psi_\jv + \Dzx\Phi_\jv \Big)  \qquad \jv\in\p\Omega_{h,r} . \label{eq:sommerfeldEigenvalueV1A}
\ea
\ese
Condition~\cref{eq:sommerfeldEigenvalueV1A} could be used to eliminate the ghost point value of either $\Phi_\jv$ or $\Psi_\jv$ in~\cref{eq:discreteEigenvalueGeneralized} 
but not both ghost points.
Thus one ghost point remains in the eigenvalue problem after eliminating constraints.

\textbf{Special case 1. Ghost points can be eliminated} Take the case of explicit time-stepping and the Sommerfeld
RBC
\ba
   \Dzt W_\jv^2 - \sigma c \Dzx W_\jv^n = 0.
\ea
which transforms to 
\ba
   \f{A-A^{-1}}{2\dt} \Phi_\jv - \sigma c \Dzx \Phi_\jv, \label{eq:SRBC-E2}
\ea
Now~\cref{eq:SRBC-E2} can be used to eliminate the ghost points for $\Phi_\jv$ in
\ba
   \f{ A - 2 + A^{-1} }{\dt^2}\Phi_\jv = L_h \Phi_\jv , \label{eq:discreteEigenvalue}
\ea
and then no ghost points are needed in the eigenvalue problem.


\textbf{Special case 2. RBC eigenvalue is consistent.}
Take the case of implicit time-stepping with full-weighting, $\beta_2=1/4$ together with the Sommerfeld RBC
\ba
   \Dpt W_\jv^n - \sigma c \Dzx( W_\jv^{n+1} + W_\jv^{n} )/2 , \qquad \jv\in\p\Omega_{h,r},
\ea
Now we have the interior eigenvalue equation 
\ba
  \lambdaTilde^2 \eqdef  - \f{4}{\dt^2} \f{A - 2 + A^{-1}}{A + 2 + A^{-1}} =  - \f{1}{(\dt/2)^2} \f{A^2 - 2A + 1}{A^2 + 2A +1 } .
\ea
and the S-RBC eigenvalue equation
\ba
  \f{1}{(\dt/2)} \f{A-1}{A+1} \Phi_\jv - \sigma c \Dzx \Phi_\jv = 0 ,  \\
  i\sigma \lambdaHat  \Phi_\jv - \sigma c \Dzx \Phi_\jv = 0 ,  
\ea
with $i\sigma \lambdaHat = \f{1}{(\dt/2)} \f{A-1}{A+1}$.
We now have $(i\sigma\lambdaHat) = -\lambdaTilde^2$ and so $\lambdaTilde=\lambdaHat$ (since we need to choose the square root to be consistent with the RBC).
{\blue For this scheme, the eigenvector $\Phi_\jv$ will correspond to the direct discretization of the eigenvalue problem in~\cref{eq:contintuousEigenvalueProblem}.
}

% To form the linear eigenvalue problem set
% \ba
%    \Chi_\jv = A{-1} \Phi_\jv
% \ea
% and~~\cref{eq:discreteEigenvalue} becomes
% \ba
%    &  A\Phi_\jv - 2 \Phi_\jv + \Chi_\jv =  \dt^2 L_h \Big( A  \alpha_2 \Phi_\jv + \beta_2 \Phi_\jv + \alpha_2 \Chi_\jv \Big), \\
%    &  \Phi_\jv = A \Chi_\jv  
% \ea
% or 


% -----------------------------------------------------------------
\clearpage
\input tex/waveHoltzResults

% -----------------------------------------------------------------
\clearpage
\input tex/eigenWaveResults

\appendix

\input tex/timeCorrections

%  \clearpage
% \bibliographystyle{elsart-num}
% \bibliography{journal-ISI,henshaw,henshawPapers}
 
\end{document}

